\chapter{Cataract Surgery}

\section*{Overview}
Cataract surgery removes a cloudy natural lens from the eye and replaces it with an artificial intraocular lens (IOL) to improve vision. The exact approach, setting, and preparation depend on the indication, anatomy, and the patient's health.

\section*{Why It Is Done}
It is performed when a cataract causes visual impairment that interferes with reading, driving, work, or daily activities and glasses no longer provide adequate improvement.

\section*{How It Is Performed}
After local anesthesia, the surgeon makes a small corneal incision, breaks up and removes the cloudy lens, and inserts a foldable IOL. The incision usually seals without stitches.

\section*{Risks and Recovery}
Infection, bleeding, inflammation, retinal problems, increased eye pressure, lens displacement, and posterior capsule clouding are possible. Use prescribed eye drops, avoid contamination or heavy strain as instructed, and attend follow-up visits.

\section*{Outcome and Follow-up}
The expected outcome depends on the reason for the procedure, the person's health, and whether additional treatment is needed. The clinician reviews the result with the patient, explains any next steps, and sets follow-up according to the findings and recovery.

\section*{When to Seek Medical Care}
Follow the procedure team's specific instructions. Seek urgent medical assessment for severe or worsening pain, uncontrolled bleeding, fever or signs of infection, trouble breathing, chest pain, fainting, new weakness or confusion, inability to urinate, or any rapidly worsening symptom. The appropriate warning signs depend on the procedure and the patient's medical condition.

\section*{References}
\begin{enumerate}
  \item MedlinePlus. Cataract Removal. U.S. National Library of Medicine; 2025.
  \item American Academy of Ophthalmology. Cataract Surgery; 2025.
\end{enumerate}
\clearpage
